\begin{description}
\item[a.] We start from the definition of $r$:
\[
r = \frac{\Delta E}{\Delta t \cdot S_{\text{stea}} \cdot \Delta\lambda}
  = \frac{\Delta E}{\Delta t \cdot 4\pi R^2 \cdot \Delta\lambda}
\quad\text{where $R$ is the radius of the star}
\]
% "S_stea" (Romanian: star surface) as in the source
\[
r = \frac{2\pi h c^2}{\lambda^5 \left( e^{hc/k\lambda T} - 1 \right)}
\]
Considering $d$ as the distance from the star to the Earth, the
definition-relation of the spectral intensity can be written as follows:
\begin{align*}
I(\lambda) &= \frac{\;\dfrac{\Delta E}{\Delta t \cdot \Delta\lambda}\;}{4\pi d^2}\\
I(\lambda) &= \frac{2\pi h c^2 R^2}
                  {d^2 \lambda^5 \left( e^{hc/\lambda k T} - 1 \right)};
\end{align*}
Particularly for each wavelength:
\[
I_1(\lambda_1) = \frac{2\pi h c^2 R^2}
  {d^2 \lambda_1^5 \left( e^{hc/\lambda_1 k T} - 1 \right)};\qquad
I_2(\lambda_2) = \frac{2\pi h c^2 R^2}
  {d^2 \lambda_2^5 \left( e^{hc/\lambda_2 k T} - 1 \right)};
\]
The ratio of the 2 above relations
\begin{equation*}
\frac{I_1(\lambda_1)}{I_2(\lambda_2)}
  = \left( \frac{\lambda_2}{\lambda_1} \right)^{\!5}
    \cdot \frac{e^{hc/\lambda_2 k T} - 1}{e^{hc/\lambda_1 k T} - 1}
\tag{1}
\end{equation*}
Represents an equation which allow to find out the temperature of star's
surface $T$ by using spectral measurements.

\item[b.] If we consider that $hc \gg \lambda k T$, then:
\[
e^{hc/\lambda_1 k T} - 1 \approx e^{hc/\lambda_1 k T}
\quad\text{and}\quad
e^{hc/\lambda_2 k T} - 1 \approx e^{hc/\lambda_2 k T}
\]
The relation (1) becomes
\[
\frac{I_1(\lambda_1)}{I_2(\lambda_2)}
  = \left( \frac{\lambda_2}{\lambda_1} \right)^{\!5}
    \cdot \frac{e^{hc/\lambda_2 k T}}{e^{hc/\lambda_1 k T}}
  = \left( \frac{\lambda_2}{\lambda_1} \right)^{\!5}
    \cdot e^{hc/\lambda_2 k T - hc/\lambda_1 k T}
\]
\[
\frac{I_1(\lambda_1)}{I_2(\lambda_2)}
  \cdot \left( \frac{\lambda_1}{\lambda_2} \right)^{\!5}
  = e^{hc/\lambda_2 k T - hc/\lambda_1 k T}
\]
\[
T = \frac{hc\,(\lambda_1 - \lambda_2)}
         {k \lambda_1 \lambda_2 \cdot
          \ln\!\left[ \dfrac{I_1(\lambda_1)}{I_2(\lambda_2)}
          \cdot \left( \dfrac{\lambda_1}{\lambda_2} \right)^{\!5} \right]}
\]

\item[c.] If $hc \ll \lambda k T$, then:
\begin{align*}
\frac{hc}{k\lambda_1 T} \ll 1;\quad
e^{hc/\lambda_1 k T} - 1 &\approx 1 + \frac{hc}{k\lambda_1 T} - 1
  = \frac{hc}{k\lambda_1 T}\\
\frac{hc}{k\lambda_2 T} \ll 1;\quad
e^{hc/\lambda_2 k T} - 1 &\approx 1 + \frac{hc}{k\lambda_2 T} - 1
  = \frac{hc}{k\lambda_2 T}
\end{align*}
The relation (1) becomes:
\[
\frac{I_1(\lambda_1)}{I_2(\lambda_2)}
  = \left( \frac{\lambda_2}{\lambda_1} \right)^{\!5}
    \cdot \frac{\;\dfrac{hc}{k\lambda_2 T}\;}{\dfrac{hc}{k\lambda_2 T}}
% sic: the source prints k lambda_2 T in both numerator and denominator;
% the denominator should be k lambda_1 T
  = \left( \frac{\lambda_2}{\lambda_1} \right)^{\!4}
\]
\[
I_1 = 2 I_2;\quad
\left( \frac{\lambda_2}{\lambda_1} \right)^{\!4} = 2;\quad
\lambda_2 = \lambda_1 \cdot \sqrt[4]{2} \approx 1{,}2 \cdot \lambda_1.
\]
\end{description}
